% Rate conversion and scale selection in the small-support regular source.
\section{Rate conversion and the local scale budget}
\label{sec:ex-lattice-trace-local-rate-selection}

This section records the quantitative selection after the Fourier-tail and
jet comparison in Section~\ref{sec:ex-lattice-trace-row-jet-estimates}.
It proves the finite-rate summation used to select the probe estimate and the
scalar domination used to choose the local chart budget. The strip, growth,
weighted-vanishing, chart, and support receipts themselves remain the inputs
identified in that section; the evidence record supplies the source crosswalk.

Fix the block side \(L\). Write
\[
 \begin{gathered}
 H(x)=\operatorname{fineMonitor}(x),\quad \chi(x)=\sqrt{H(x)},\quad
 \beta(x)=\beta_F(x),\\
 R(x)=\operatorname{radiusD}(x),\quad
 \rho(x)=\operatorname{reach}(x),\quad b(x)=\operatorname{radiusB}(x).
 \end{gathered}
\]
For fixed \(L\), these are
\[
 \begin{gathered}
 \beta(x)=x+c_L,\quad b(x)=\lceil(\log x)^2\rceil,\\
 R(x)=b(x)(4b(x)+2)+2b(x),\quad \rho(x)=45(1+b(x)),\\
 \operatorname{depthW}_L(R)=L(R+2),
 \end{gathered}
\]
where \(c_L\) is fixed. Here \(H(x)=\exp((\log\log\beta(x))^2)\), so \(H\ge1\),
\(H\to\infty\), and \(\chi\) is the monitor factor in the priced
conclusion. For a function \(f\), write \(f\in\mathcal R_a\) if for
all \(v>a\) one has eventually \(|f(x)|\le\beta(x)^v\). This is the
source's strict-power rate convention. Finite sums preserve a common rate,
products add rates, and \(\chi\in\mathcal R_0\).

\begin{lemma}[The exponential tail pays the finite probe]
\label{lem:ex-lattice-tail-pays-probe}
Let \(d\in\mathbb N\), \(\sigma>0\), and let \(K,M\) and the priced
local derivative bound \(\Phi\) be fixed nonnegative constants. Let
\(\delta_k(x)\) and \(C_1(x)\) be the quantities of
(\ref{eq:ex-lattice-local-delta})--(\ref{eq:ex-lattice-local-C1}), with
\(b=\rho(x)\) and \(m=R(x)\). Set
\[
 E(x)=\sum_{k=0}^{d}\delta_k(x),\qquad
 Z_d(x)=\sum_{j=0}^{d}\#\operatorname{Part}^{\mathrm{ord}}(j)\,2^j
       E(x)C_1(x)^j\,\Phi.
 \tag{12}\label{eq:ex-lattice-tail-probe-sum}
\]
Then \(Z_d\in\mathcal R_{-1}\), and consequently
\[
                 \chi(x)Z_d(x)\longrightarrow0.
 \tag{13}\label{eq:ex-lattice-tail-probe-price}
\]
In particular, for every \(y>0\), eventually
\(Z_d(x)\le y/\chi(x)\). If the local receipts of
Proposition~\ref{prop:ex-lattice-strip-to-probe} hold and
\(E(x)\le C_1(x)\), every derivative order \(j\le d\) of the selected
probe is bounded by \(Z_d(x)w\), where \(w\) is the derivative bound of
the observable appearing in (\ref{eq:ex-lattice-jet-price}) and
the price factor \(\Phi\) has already been included in \(Z_d\).
\end{lemma}
\begin{proof}
By the polynomial geometry receipts, \(R\) and the lift-size polynomial
\((1+4L(R+2))^{dk}\) have rate zero. The width cost and all other factors
in \(\delta_k\), apart from \(\exp(-\sigma\rho(x)/2)\), are fixed constants
or rate-zero polynomial factors. The reach cutoff has every negative rate:
for each \(a>0\), eventually \(\exp(-\sigma\rho(x)/2)\le\beta^{-a}\).
Thus each \(\delta_k\), and their finite sum \(E\), have rate \(-1\)
(the stronger arbitrary negative rate is more than needed). The finite sum
in \(C_1\) has rate zero. By the product and finite-sum rules,
\(Z_d\in\mathcal R_{-1}\).

Since \(\chi\in\mathcal R_0\), their product has rate \(-1\). The
strict-power definition then implies convergence to zero: choose any
exponent \(v\) strictly between \(-1\) and \(0\); eventually
\(|\chi Z_d|\le\beta^v\), and \(\beta\to\infty\). This proves (13)
and the stated eventual price bound. The final assertion follows by
applying Proposition~\ref{prop:ex-lattice-strip-to-probe} at each order and
bounding its finite ordered-partition sum by (12).
\end{proof}

\paragraph{Actual finite-rate receipts.}
In the construction the polynomial direction-growth constants are
\[
 h_k(x)=4K M\,C_{dk}(\sigma)\bigl(1+4L(R(x)+2)\bigr)^{dk},\qquad
 C_1(x)=1+\sum_{k=0}^d h_k(x).
\]
The error at order \(k\) is
\[
 \delta_k(x)=8KM e^{-\sigma \rho(x)/2}
   (1+4R(x))^{dk}C_{dk}(\sigma/2).
\]
These formulas retain the source's factor four from the four edge
components and the additional factor two from the two input arrays. The
coefficient \(E(x)\) is required to be at most \(C_1(x)\) before applying
the finite multilinear comparison. This is a separate eventual estimate;
it follows because \(E\in\mathcal R_{-1}\) and \(C_1\ge1\).

\begin{proposition}[The selected local budget]
\label{prop:ex-lattice-selected-local-budget}
Fix natural numbers \(n,p,r\) and a fixed positive integer block side \(L\). Let
\[
 C.C\ge1,\quad C.M\ge1,\quad D.Q\ge1,\quad D.K\ge1,\quad D.Z\ge1,
 \quad D.\varepsilon>0
\]
be fixed real constants; in the selected source chart they are scalar
entries of the chart data \(C\) and of the record \(D\), and the
dotted notation is kept for comparison with the source. These ranges are
the only meaning of ``admissible'' below, and no other property of \(C\)
and \(D\) is used except the nonnegative linear clock introduced below.
Put
\[
 P_n=\sum_{j=0}^{n}\#\operatorname{Part}^{\mathrm{ord}}(j),\qquad
 \operatorname{clinicFee}(n,Q)=2Q(1+P_n),
\]
and let \(\mathcal C_L\) be the finite index set
\(\mathrm{Site}_L\times\mathrm{Axis}\times I_G\), where \(\mathrm{Site}_L\)
is the set of sites of the periodic lattice of side \(L\) and \(I_G\) is
the finite set of colour coordinates, of cardinality \(\dim\mathfrak g\).
For
\(x\) large, define the covering and response quantities
\begin{align*}
 P_*(x)&=(2\,\operatorname{depthW}_L(R(x))+1)^4,\\
 B_*(x)&=(r+48)(1+b(x))(1+P_*(x)),\\
 A_*(x)&=\operatorname{boost}_{C,D,L}\,B_*(x)^{\,8+2p+2pn}
                 +\frac{16}{D.\varepsilon},\\
 S_*(x)&=\operatorname{clinicFee}(n,D.Q)
              \bigl(2C.C(1+R(x))^p\bigr)^2,\\
 U_*(x)&=(A_*(x)+S_*(x))\,\operatorname{fiery}(x)
                 +\frac{2A_*(x)+2}{H(x)},\\
 K_*(x)&=1+\operatorname{guests}_{G,D.K}(n+1,L)
                      B_*(x)^{\,4(1+4(n+1))}.
\end{align*}
These abbreviations use the exact source coefficients
\[
\begin{aligned}
 s_L&=1+16\,\#\mathcal C_L,\\
 b_0&=(\operatorname{clinicFee}(n,D.Q)+C.M)(2C.C)^2,\\
 \operatorname{boost}_{C,D,L}&=b_0+(b_0+1)L^3
                     +2n s_L^2(D.Z+1)D.Z,\\
 \operatorname{guests}_{G,D.K}(q,L)&=2s_L^{1+4q}
                     \operatorname{inflate}_{4q}(D.K^{-1}/8),\\
 \operatorname{meter}(h,D.K,m)&=2(1+h)^{m+1}
                     \operatorname{inflate}_{m}(D.K^{-1}/8).
\end{aligned}
\tag{14a}\label{eq:ex-lattice-scale-coefficients}
\]
Here \(h\) is the number of heads, the subscript of
\(\operatorname{guests}\) records its dependence on the gauge group
(through \(I_G\)) and on the scalar \(D.K\), and
\(\operatorname{inflate}_m(v)=(1+m/v)^m\) is the Gaussian moment
inflation factor. The scalar cost terms below
are nonnegative by the stated ranges of the constants. The source clock is
a fixed function of the form
\(\operatorname{clocks}(t)=c_{\mathrm{clock}}t\) with
\(c_{\mathrm{clock}}=\operatorname{clocks}(1)\ge0\).
Let
\(\beta\ge\beta(x)/2\). Suppose the support and budget satisfy
\[
 \begin{gathered}
 B.\mathrm{threshold}=\operatorname{fiery}(x),\quad
 0\le P.\mathrm{threshold}\le2B.\mathrm{threshold},\\
 B.B\le2b(x),\quad P.\mathrm{power}=p,\quad
 \operatorname{coverNumber}(B.B,X)\le P_*(x).
 \end{gathered}
 \tag{14}\label{eq:ex-lattice-budget-support}
\]
Put \(F_0=B.\mathrm{threshold}\), \(u=\beta^{-1/2}\),
\(v_\rho=C.C\,P.\mathrm{threshold}(1+\rho(x))^p\),
\(v_R=C.C\,P.\mathrm{threshold}(1+R(x))^p\),
\(h=\operatorname{heads}(B,X)\), and
\(\operatorname{cup}=D.Z(1+\operatorname{scope}(B,X,r,x))^{pn}\).
Define the displayed scalar cost by
\[
\begin{aligned}
 F(x)={}&(\operatorname{clinicFee}(n,D.Q)+C.M)v_\rho^2
       (2+\operatorname{breadth}(B,X))^{pn}/F_0\\
 &+\frac{\operatorname{clinicFee}(n,D.Q)v_\rho^2+u}{F_0}L^3\\
 &+n(2F_0)h^2(\operatorname{cup}+1)\operatorname{cup}\\
 &+\frac{16}{D.\varepsilon}(F_0+u/F_0)\\
 &+\frac{\operatorname{clinicFee}(n,D.Q)v_R^2+u}{F_0}.
\end{aligned}
\tag{14b}\label{eq:ex-lattice-explicit-local-cost}
\]
This is the response-plus-paint cost used by the local chart construction.
It satisfies,
eventually in \(x\),
\[
 0\le F(x)\le U_*(x),\qquad
 \operatorname{meter}(\operatorname{heads}(B,X),D.K,4(n+1))\le K_*(x),
 \tag{15}\label{eq:ex-lattice-budget-majorants}
\]
and
\[
 \chi(x)U_*(x)K_*(x)\longrightarrow0.
 \tag{16}\label{eq:ex-lattice-budget-small}
\]
Consequently, for every \(y>0\), eventually and uniformly in all data
satisfying (14),
\[
 F(x)\le1,\qquad
 \operatorname{clocks}(t)F(x)\operatorname{meter}(\operatorname{heads}(B,X),D.K,4(n+1))
       \le \frac{y}{\chi(x)}t\quad(t\ge0).
 \tag{17}\label{eq:ex-lattice-budget-output}
\]
This is the scalar estimate consumed by the local regular-source branch.
\end{proposition}
\begin{proof}
We give the domination and limiting arguments separately. First, the cover
formula gives
\[
 1+\rho(x)\le B_*(x),\quad
 2+\operatorname{breadth}(B,X)\le B_*(x),\quad
 1+\operatorname{scope}(B,X,r,x)\le B_*(x).
 \tag{18}\label{eq:ex-lattice-budget-cover}
\]
Write \(a=b(x)\) and \(q=\operatorname{coverNumber}(B.B,X)\). The
definitions give
\[
 \operatorname{breadth}(B,X)=q(4B.B+2)+2B.B+8
       \le q(8a+2)+4a+8,
 \quad \operatorname{scope}=r a+\operatorname{breadth}+1.
\]
With \(S_X=(r+48)(1+a)(1+q)\), direct expansion, using \(a,q,r\ge0\),
gives \(1+45(1+a)\le S_X\),
\(2+q(8a+2)+4a+8\le S_X\), and
\(r a+q(8a+2)+4a+10\le S_X\). These are (18) with \(S_X\) in place
of \(B_*\). Since \(q\le P_*\) and the bulk formula is increasing in
its cover-count argument, \(S_X\le B_*\); hence (18) follows.

Put \(F_0=B.\mathrm{threshold}\), \(u=gStep(\beta)=\beta^{-1/2}\),
and \(S_X=\operatorname{bulk}(\operatorname{coverNumber}(B.B,X),r,x)\).
Since \(F_0=\beta(x)^{-1/2}H(x)\), (14) yields
\[
 \frac{u}{F_0}
 =\frac{1}{H(x)}\sqrt{\frac{\beta(x)}{\beta}}
 \le\frac{\sqrt2}{H(x)}\le\frac{2}{H(x)}.
 \tag{19}\label{eq:ex-lattice-u-over-F}
\]
The cover inequalities (18) and \(P.\mathrm{threshold}\le2F_0\)
give
\[
 \operatorname{loom}(P,\rho(x),C.C)\le
             2C.C F_0 S_X^p.
 \tag{20}\label{eq:ex-lattice-loom-bound}
\]
Put \(t=F_0\) and \(v=\operatorname{loom}(P,\rho(x),C.C)\). The loom is
\(C.C\,P.\mathrm{threshold}(1+\rho(x))^p\); (18) and
\(P.\mathrm{threshold}\le2t\) give \(v\le2C.CtS_X^p\). The segment
term is \((\operatorname{clinicFee}+C.M)v^2\operatorname{heft}/t\),
where \(\operatorname{heft}=(2+\operatorname{breadth})^{pn}\le
S_X^{pn}\). Thus, with
\(b_0=(\operatorname{clinicFee}+C.M)(2C.C)^2\),
\[
 \operatorname{segment}\le b_0 t S_X^{2p+pn}
       \le b_0\zeta(B,u)S_X^{8+2p+2pn}.
\]
The tranche term is
\(((\operatorname{clinicFee}v^2+u)/t)L^3\). Its first numerator part
is bounded by \(b_0tS_X^{2p}\); its second is \(u/t\). Since
\(t\le\zeta(B,u)=t+u/t\), \(S_X\ge1\), and
\(2p\le8+2p+2pn\), this gives
\(\operatorname{tranche}\le(b_0+1)L^3\zeta(B,u)
S_X^{8+2p+2pn}\).

For the scratch term, its exact factors are
\(n(2t)\operatorname{heads}^2(\operatorname{cup}+1)\operatorname{cup}\),
where \(\operatorname{cup}=D.Z(1+\operatorname{scope})^{pn}\).
The head count is
\((2\operatorname{scope}+1)^4\#\mathcal C_L\). From (18),
\(1+\operatorname{heads}\le s_LS_X^4\), with
\(s_L=1+16\#\mathcal C_L\).
Also \(\operatorname{cup}\le D.ZS_X^{pn}\) and
\(\operatorname{cup}+1\le(D.Z+1)S_X^{pn}\). Replacing \(t\) by
\(\zeta(B,u)\) therefore gives
\[
 \operatorname{scratch}\le
 2n s_L^2(D.Z+1)D.Z\,\zeta(B,u)S_X^{8+2pn}.
\]
Adding these three bounds gives the coefficient
\(\operatorname{boost}=b_0+(b_0+1)L^3+2n s_L^2(D.Z+1)D.Z\).
Since \(\operatorname{opal}=\operatorname{rust}+(16/D.\varepsilon)\zeta\),
and \(S_X\ge1\),
\[
 \operatorname{opal}\le A_*(x)(F_0+u/F_0).
 \tag{21}\label{eq:ex-lattice-opal-bound}
\]
The paint term is evaluated with the loom at \(R(x)\), not at the reach.
Its direction is
\(v_R=C.C\,P.\mathrm{threshold}(1+R(x))^p\le
2C.CF_0(1+R(x))^p\). Squaring this inequality and dividing by \(F_0>0\)
gives \(\operatorname{clinicFee}(n,D.Q)v_R^2/F_0\le S_*(x)F_0\).
The remaining paint term is \(u/F_0\). Thus the paint contribution is bounded by
\[
 \operatorname{paintW}\le S_*(x)F_0+u/F_0.
 \tag{22}\label{eq:ex-lattice-paint-bound}
\]
Equations (19)--(22) give \(F\le U_*\).

The head-count moment estimate follows by inserting
\(1+\operatorname{heads}\le s_LS_X^4\) into the explicit moment formula
(14a): for \(m=4(n+1)\),
\[
 \operatorname{meter}(\operatorname{heads}(B,X),D.K,4(n+1))
 \le \operatorname{guests}_{G,D.K}(n+1,L)S_X^{4(1+4(n+1))}
 \le K_*(x).
 \tag{23}\label{eq:ex-lattice-head-moment}
\]

For the limit, \(B_*,A_*,S_*,K_*\) are fixed polynomials in
\(1+\operatorname{radiusB}(x)\), hence have rate zero and are subpower
relative to \(H\). Indeed \(\operatorname{radiusB}(x)\le
1+(\log x)^2\), while \(\beta(x)=x+O(1)\) for fixed \(L\), and
\(\log H=(\log\log\beta)^2\); the logarithm of any such fixed
polynomial is \(O(\log\log\beta)=o(\log H)\). Thus for each \(\eta>0\),
eventually each nonnegative polynomial factor \(Q\) here satisfies
\(Q\le H^\eta\). The first term in \(\chi U_*K_*\) has rate
\(-1/2\), since \(\operatorname{fiery}=H/\sqrt{\beta}\) and all other
factors have rate zero. For the second term choose the subpower bounds
\(A_*+1\le2H^{1/8}\), \(K_*\le H^{1/8}\); then
\[
 \chi\,\frac{(2A_*+2)K_*}{H}
 \le4H^{1/2+1/8+1/8-1}=4H^{-1/4}\longrightarrow0.
\]
This proves (16). Also \(H\ge1\), \(K_*\ge1\), and \(F\le U_*\), so
\(F\le\chi U_*K_*\); (16) makes this at most one eventually. Finally
\(\operatorname{clocks}(t)=t\operatorname{clocks}(1)\) for \(t\ge0\).
Choose \(d_*=1+\operatorname{clocks}(1)\). From (16), eventually
\(d_*\chi U_*K_*\le y\), and (15) proves the second inequality in (17).
\end{proof}

The bound is conditional only on the selected local chart and support
inequalities (14); the source crosswalk identifies the already-proved
support-cover, rust, and Gaussian-moment estimates whose algebra is expanded
above. It does not produce the separate strip, field-growth, weighted
vanishing, same-history, or healthy-chart receipts required to enter this
local branch.
